
\begin{appendices}

\section{The safe linear e-process}\label{app:safelinear}

We derive~\eqref{eq:safelinear} and prove Proposition~\ref{prop:safelinear}.

\subsection{Setup}

Write $y \in \Real^n$ for the stacked observations of $X_i$, $x \in \Real^n$
for those of $X_j$, and $W \in \Real^{n \times q}$ for the design matrix of
the conditioning block $X_S$ together with an intercept, $q = |S| + 1$. The
model is
\begin{equation}\label{eq:appmodel}
y \;=\; W\alpha \;+\; \beta x \;+\; \varepsilon,
\qquad \varepsilon \sim N(0, \sigma^2 I_n),
\end{equation}
with $H_0 : \beta = 0$ and $H_1 : \beta \neq 0$, and nuisance parameters
$(\alpha, \sigma) \in \Real^q \times (0,\infty)$ common to both.

Let $M = I_n - W(W^\top W)^{-1} W^\top$ be the residual projector onto the
orthogonal complement of the column space of $W$, and write
$\tilde y = My$, $\tilde x = Mx$ for the residualised vectors. Set
\begin{equation*}
S_{yy} = \tilde x^\top \tilde x, \qquad
S_{xy} = \tilde x^\top \tilde y, \qquad
S_{xx} = \tilde y^\top \tilde y,
\end{equation*}
following the naming convention of the implementation, in which the first
subscript refers to the \emph{response} of the regression being examined. The
squared sample partial correlation is
$\hat r^2 = S_{xy}^2 / (S_{xx} S_{yy})$, and the least-squares estimate of
$\beta$ is $\hat\beta = S_{xy}/S_{yy}$.

\subsection{Invariance and the right-Haar prior}

The null model~\eqref{eq:appmodel} with $\beta = 0$ is invariant under the
group $G$ of transformations
\begin{equation*}
g_{c,\lambda} : \quad y \mapsto \lambda y + Wc,
\qquad c \in \Real^q, \ \lambda > 0,
\end{equation*}
acting on the parameters by $\alpha \mapsto \lambda\alpha + c$,
$\sigma \mapsto \lambda\sigma$, and---under the alternative---$\beta \mapsto
\lambda\beta$ once $x$ is transformed correspondingly. The right-Haar measure
for this group is $\pi(\alpha,\sigma)\,d\alpha\,d\sigma \propto
\sigma^{-1}\,d\alpha\,d\sigma$.

The relevant fact \citep{berger1998bayes,hendriksen2021optional} is that when
the same right-Haar prior is placed on a nuisance group shared by both
hypotheses, the resulting Bayes factor equals the Bayes factor based on the
maximal invariant, whose distribution under the null does not depend on the
nuisance parameters. Consequently the Bayes factor is a \emph{test
martingale}: its expectation under any null distribution, conditional on the
past, is one, whatever the nuisance parameters happen to be. This is the
property that fails for a generic Bayes factor and that makes optional
stopping legitimate here.

\subsection{Integrating the marginal likelihoods}

Under $H_0$, integrating~\eqref{eq:appmodel} against $\pi(\alpha,\sigma)
\propto \sigma^{-1}$ gives
\begin{equation}\label{eq:m0}
m_0(y) \;=\; \int_0^\infty \int_{\Real^q}
  (2\pi\sigma^2)^{-n/2}
  \exp\!\left\{-\frac{\|y - W\alpha\|^2}{2\sigma^2}\right\}
  \frac{d\alpha \, d\sigma}{\sigma}
\;=\; C_{n,q} \, \bigl(S_{xx}\bigr)^{-(n-q)/2},
\end{equation}
where $C_{n,q} = \tfrac12 \pi^{-(n-q)/2} \Gamma\!\bigl(\tfrac{n-q}{2}\bigr)
|W^\top W|^{-1/2}$ collects the terms that do not depend on the data through
the residuals. The $\alpha$-integral is Gaussian and contributes
$|W^\top W|^{-1/2}$ together with the replacement of $\|y-W\alpha\|^2$ by
$\|\tilde y\|^2 = S_{xx}$; the $\sigma$-integral is then a gamma integral.

Under $H_1$ we add $\beta \mid \sigma \sim N(0, \sigma^2 g)$, independent of
$\alpha$. Completing the square in $\beta$ against the residualised design
gives a Gaussian integral with precision $\sigma^{-2}(S_{yy} + g^{-1})$ and
mean $S_{xy}/(S_{yy} + g^{-1})$, so
\begin{align}
m_1(y)
&= C_{n,q} \, \bigl(1 + gS_{yy}\bigr)^{-1/2}
   \left( S_{xx} - \frac{S_{xy}^2}{S_{yy} + g^{-1}} \right)^{-(n-q)/2}
   \notag \\[2pt]
&= C_{n,q} \, (1+G)^{-1/2} \,
   S_{xx}^{-(n-q)/2} \bigl(1 - \kappa \hat r^2\bigr)^{-(n-q)/2},
   \label{eq:m1}
\end{align}
where $G = gS_{yy}$ and $\kappa = G/(1+G)$, using
$\dfrac{S_{xy}^2}{S_{xx}(S_{yy}+g^{-1})}
 = \dfrac{S_{xy}^2}{S_{xx}S_{yy}} \cdot \dfrac{S_{yy}}{S_{yy}+g^{-1}}
 = \hat r^2 \kappa$.

Dividing~\eqref{eq:m1} by~\eqref{eq:m0}, the constants and the
$S_{xx}^{-(n-q)/2}$ factors cancel:
\begin{equation}\label{eq:bf}
E_n \;=\; \frac{m_1(y)}{m_0(y)}
\;=\; (1+G)^{-1/2} \bigl(1 - \kappa \hat r_n^2\bigr)^{-(n-q)/2},
\end{equation}
which is~\eqref{eq:safelinear} with $p = q = |S|+1$. Proposition~%
\ref{prop:safelinear} follows from the invariance argument above:
$(E_n)$ is a ratio of marginal likelihoods under a shared right-Haar prior, so
it is a test martingale under every null distribution, and by optional
stopping an e-process.

\subsection{Two implementation notes}

\paragraph{Fixing $g$.} $G = gS_{yy}$ grows with $n$ because $S_{yy}$ does.
The martingale property requires $g$ to be $\Fcal_0$-measurable; setting it
from the running data---for instance to hold $G$ constant---breaks it. We
choose $g$ from a warm-up prefix that is discarded, so it is a constant as far
as every e-process is concerned.

\paragraph{Streaming.} All of $S_{xx}, S_{xy}, S_{yy}$ are Schur complements
of the running Gram matrix of the full observation vector. Maintaining that
matrix costs $O(d^2)$ per observation and serves every $(i,j,S)$ triple, so
adding hypotheses costs computation at query time but no additional state.

\section{Inverting the e-process into a confidence sequence}\label{app:cs}

Fix $\beta_0$ and consider the null $\beta = \beta_0$. Replacing $y$ by
$y - \beta_0 x$ in~\eqref{eq:bf} gives the e-process for that null; the
confidence sequence at level $1-\alpha$ is the set of $\beta_0$ not rejected,
\begin{equation*}
\Bigl\{ \beta_0 :
  (1+G)^{-1/2}\bigl(1 - \kappa \hat r^2(\beta_0)\bigr)^{-(n-q)/2}
  < 1/\alpha \Bigr\},
\end{equation*}
where $\hat r^2(\beta_0)$ is the squared partial correlation computed from
$y - \beta_0 x$. Taking logarithms, the condition reads
\begin{equation*}
-\tfrac12 \log(1+G) - \tfrac{n-q}{2}\log\bigl(1 - \kappa\hat r^2(\beta_0)\bigr)
  < -\log\alpha ,
\end{equation*}
that is $\hat r^2(\beta_0) < \tau$ with
\begin{equation}\label{eq:tau}
\tau \;=\; \frac{1}{\kappa}
  \left[ 1 - \exp\!\left( \frac{2}{n-q}
    \Bigl( \log\alpha - \tfrac12\log(1+G) \Bigr) \right) \right]
\;=\; \frac{1}{\kappa}
  \left[ 1 - \left(\frac{\alpha}{\sqrt{1+G}}\right)^{2/(n-q)} \right].
\end{equation}
Writing $S_{xy}(\beta_0) = S_{xy} - \beta_0 S_{yy}$ and
$S_{xx}(\beta_0) = S_{xx} - 2\beta_0 S_{xy} + \beta_0^2 S_{yy}$, the condition
$S_{xy}(\beta_0)^2 < \tau\, S_{xx}(\beta_0) S_{yy}$ is a quadratic inequality
in $\beta_0$ with leading coefficient $(1-\tau)S_{yy}^2 > 0$ whenever
$\tau < 1$, and solving it gives
\begin{equation}\label{eq:csapp}
\hat\beta \;\pm\; \frac{1}{S_{yy}}
  \sqrt{\frac{\tau\bigl(S_{xx}S_{yy} - S_{xy}^2\bigr)}{1-\tau}},
\qquad \hat\beta = S_{xy}/S_{yy},
\end{equation}
which is~\eqref{eq:cs}. When $\tau \le 0$---which happens at small $n$, since
$(\alpha/\sqrt{1+G})^{2/(n-q)} \to 1$ from above as $n \downarrow q$---the
inequality has no solution and the interval is the whole line.

\begin{remark}[A sign that mattered]
The bracketed term in~\eqref{eq:tau} carries $-\tfrac12\log(1+G)$, not
$+\tfrac12\log(1+G)$. With the sign reversed, the formula still produces
intervals of a plausible shape, still shrinks at the right rate, and still
passes every Type-I error simulation of the underlying \emph{test}, because
the test is unaffected. It produces intervals roughly half the correct width.
We caught it only by comparing~\eqref{eq:csapp} against a brute-force
numerical inversion of~\eqref{eq:bf} over a grid of $\beta_0$, which now runs
as a unit test. Closed-form inversions of e-processes deserve this check as a
matter of routine.
\end{remark}

\section{The frozen conditioning block}\label{app:conditioning}

Proposition~\ref{prop:direction} is stated for a fixed null. In
Algorithm~\ref{alg:certcd} the conditioning block $Z_{ij}$ is selected at the
random time $T$ at which the pair's adjacency is certified, so the hypothesis
being tested is itself chosen using data. We record why the guarantee survives.

Let $T$ be the certification time of pair $(i,j)$; it is a stopping time with
respect to $(\Fcal_t)$, since certification is triggered by the running
maximum of $A_t(i,j)$, which is adapted. Let $Z_{ij}$ be $\Fcal_T$-measurable
(it is: it is a function of the certified statuses at time $T$). The direction
process is built only from observations $o_{T+1}, o_{T+2}, \dots$.

Condition on $\Fcal_T$. Because the observations are i.i.d.\ and $T$ is a
stopping time, the post-$T$ observations are, conditionally on $\Fcal_T$,
i.i.d.\ from $P$; and $Z_{ij}$ is now a fixed set. The proof of
Proposition~\ref{prop:direction} therefore applies verbatim to the
conditional law, giving
$\Exp_P[E_{T+\tau} \mid \Fcal_T] \le 1$ for every stopping time $\tau$ of the
post-$T$ filtration, on the event $\{T < \infty\}$. Taking expectations,
$\Exp_P[E_{T+\tau} \mathbf{1}\{T<\infty\}] \le \Prob(T<\infty) \le 1$, and
Ville's inequality applies to the post-$T$ process. Defining $E_t = 1$ for
$t \le T$ extends this to a process on the original time index without
affecting the bound.

Two remarks. First, this is why the implementation must discard all
observations before $T$: reusing them would break the conditional
independence, and the constant they contribute is exactly the failure mode of
Remark~\ref{rem:samesupport}. Second, the argument is what licenses the
graph-dependent execution order in Theorem~\ref{thm:main}: the algorithm may
decide \emph{which} direction hypotheses to examine using the data, because
the union bound runs over the full fixed family $\{(i,j)\}$ regardless, and
each member of that family has a valid e-process whether or not the algorithm
chooses to evaluate it.

\section{Additional experimental detail}\label{app:extra}

\begin{table}[h]
\centering
\caption{Realised crossing rates of the per-set primitive under six noise
laws, $2000$ replications each, monitored to $n = 800$ after every tenth
observation. Every entry is at or below its nominal level; the largest
realised-to-nominal ratio in the table is $0.95$ (shifted exponential,
$\alpha = 0.01$) and the median is $0.49$.}%
\label{tab:primitive}
\begin{tabular}{@{}lcccc@{}}
\toprule
noise law & $\alpha=0.20$ & $\alpha=0.10$ & $\alpha=0.05$ & $\alpha=0.01$ \\
\midrule
Gaussian            & $0.091$ & $0.046$ & $0.023$ & $0.003$ \\
Laplace             & $0.092$ & $0.047$ & $0.024$ & $0.003$ \\
uniform             & $0.107$ & $0.052$ & $0.031$ & $0.006$ \\
Student-$t_5$       & $0.093$ & $0.044$ & $0.020$ & $0.004$ \\
shifted exponential & $0.101$ & $0.051$ & $0.028$ & $0.010$ \\
two-component mixture & $0.097$ & $0.049$ & $0.028$ & $0.006$ \\
\bottomrule
\end{tabular}
\end{table}

\paragraph{Hyperparameters.} Unless stated otherwise: warm-up $m = 50$
observations (discarded), prior standard deviation $0.5$ after
standardisation---so $g = 0.25$ in the parameterisation of
Eq.~\eqref{eq:safelinear}, and $g = 0.16$ in the primitive-calibration
experiment, which uses $0.4$---
$k = 2$, adjacency share $\alpha_A/\alpha = 0.5$, batch size $250$ in the
graph experiments and $100$ in the skeleton experiments, and a minimum of $80$
post-certification observations before a direction process is first evaluated.
The equivalence-region width for SPRINT-CD is $\delta = 0.15$ on the
standardised scale, matching the $\delta$ used to define the strata.

\paragraph{Reproduction.} \texttt{python experiments/run\_all.py} reproduces
every figure and table; individual scripts accept \texttt{--quick} for a
low-replication smoke test, which writes to \texttt{results/quick/} so that it
cannot overwrite the full-precision artefacts the paper cites. Seeds are fixed
in each script.

\end{appendices}

\bibliography{refs}

\end{document}
